\Needspace{6\baselineskip}
\begin{center}\hypertarget{native-266-article192}{}\pdfbookmark[2]{192.}{article-192}192.\end{center}
Methodus in art. praec. exposita praecipue determinationi \emph{primae} orbitae
perturbationes implicantis accommodata est: quamprimum vero tum elementa
media elliptica tum aequationes perturbationum proxime iam sunt cognitae,
determinatio exactissima adiumento observationum quam plurimarum commodissime
per methodum art. 187 absolvetur, quae hic explicatione peculiari
opus non habebit. Quodsi hic observationum praestantissimarum copia satis
magna est, magnumque temporis intervallum complectitur, haec methodus in
pluribus casibus simul determinationi exactiori massarum planetarum perturbantium,
saltem maiorum, inservire poterit. Scilicet, si massa cuiusdam planetae
perturbantis in calculo perturbationum supposita nondum satis certa
videtur, introducetur, praeter sex incognitas a correctionibus elementorum pendentes,
adhuc alia $\mu$, statuendo rationem massae correctae ad massam suppositam
ut $1+\mu$ ad $1$; supponere tunc licebit, perturbationes ipsas in eadem
ratione mutari, unde manifesto in singulis positionibus calculatis terminus
novus linearis ipsam $\mu$ continens producetur, cuius evolutio nulli difficultati
obnoxia erit. Comparatio positionum calculatarum cum observatis secundum
\SourceLeaf{267}principia supra exposita, simul cum correctionibus elementorum etiam correctionem
$\mu$ suppeditabit. Quinadeo hoc modo massae \emph{plurium} planetarum exactius
determinari poterunt, qui quidem perturbationes satis considerabiles
exercent. Nullum dubium est, quin motus planetarum novorum, praesertim
Palladis et Iunonis, qui tantas a Iove perturbationes patiuntur, post aliquot
decennia hoc modo determinationem exactissimam massae Iovis allaturi sint:
quinadeo forsan ipsam massam unius alteriusve horum planetarum novorum
ex perturbationibus, quas in reliquos exercet, aliquando cognoscere licebit.
\begin{center}\rule{3cm}{.4pt}\end{center}
